\documentclass[../main]{subfiles}
 
\begin{document}

\chapter{A categorical lens}


\section{A category}

A \emph{category} \(\Cat{C}\) consists of a collection of \emph{objects} \(\Ob{C}\), and for each pair of objects \(A,B \in \Ob{C}\), a collection of \emph{morphisms} \(\Hom{C}(A,B)\), also called arrows, endowed with the following properties:

\begin{enumerate}
\item given \(A,B,C \in \Ob{C}\), there exists an operation
  \[
    \circ : \Hom{C}(A,B) \times \Hom{C}(B,C) \to \Hom{C}(A,C)
  \]
  called categorical composition (which is not necessarily the composition of functions). If \(f \in \Hom{C}(A,B)\), we may write \(f \colon A \to B\), and given \(f : A \to B\) and \(g : B \to C\), we write \(g \circ f := {\circ}(f,g)\). The composition function must be associative; that is, if \(h \in \Hom{C}(C,D)\), then \((h \circ g) \circ f = h \circ (g \circ f)\);
\item for every \(A \in \Ob{C}\), there exists an identity morphism \(1_A \in \Hom{C}(A,A)\), and for all \(f,g \in \Hom{C}\) with \(\dom{f} = \cod{g} = A\), it holds that \(f \circ 1_A = f\) and \(1_A \circ g = g\).
\end{enumerate}
   
\section{Functors}
\section{Natural transformations}
\section{Limits and colimits}\label{sec:limits-colimits}


\end{document}
